\begin{table}[h]
\centering
\caption{Function calls for hierarchical retrieval}
\vspace{-8pt}
\label{tab:function_calls}
\scriptsize
\begin{tabular}{l|l|p{7.5cm}}
\hline
\textbf{Name} & \textbf{Arguments} & \textbf{Description} \\
\hline
list\_components & None &
Return all components with brief descriptions. \\
\hline
list\_packages\_under\_component & component\_name &
Return all packages under the specified component. \\
\hline
list\_entries\_under\_package & package\_name &
Return all thread entry points defined in the package. \\
\hline
list\_method\_under\_class & class\_name &
Return all methods under the specified class. \\
\hline
get\_entry\_summary & entry\_name &
Return concurrency summary of the thread entry and its interacting entries. \\
\hline
get\_method\_summary & method\_name &
Return concurrency-related information of the method. \\
\hline
get\_interaction\_dsl & entry\_pair &
Return interaction DSL for the given pair of thread entries. \\
\hline
get\_code\_snippet & method\_name &
Return code snippet of the specified method. \\
\hline
exit & None &
Terminate the retrieval process. \\
\hline
\end{tabular}
\end{table}
\vspace{-8pt}